\documentclass[11pt]{article}

% --- Packages ---
\usepackage[utf8]{inputenc}
\usepackage[margin=1in]{geometry} % Better margins
\usepackage{amsmath, amssymb, amsthm} % Essential math tools
\usepackage{hyperref} % Clickable links/TOC
\usepackage{xcolor} % Custom colors
\usepackage{listings} % For code snippets
\usepackage{enumitem} % Better control over lists
\usepackage{ulem} % For underlining
\usepackage{enumitem} % to enumerate important items with bolded numbers
\usepackage{amsmath} % to highlight important equations with a box
\usepackage{empheq}

% --- Custom Math Environments ---
\newtheorem{theorem}{Theorem}[section]
\theoremstyle{definition}
\newtheorem{definition}{Definition}[section]
\newtheorem{example}{Example}[section]

% --- Code Snippet Styling ---
\lstset{
    basicstyle=\ttfamily\small,
    breaklines=true,
    keywordstyle=\color{blue},
    commentstyle=\color{gray},
    frame=single,
    numbers=left,
    numberstyle=\tiny\color{gray}
}

\setlength{\parindent}{0pt}      % Sets indentation to zero
\setlength{\parskip}{1em}        % Adds a gap between paragraphs

% --- Document Info ---
\title{CS 109: Intro to Probability \\ \large Lecture 2: Conditional Probability and Bayes}
\author{Chris Laganiere}
\date{\today}

\begin{document}

\maketitle
\tableofcontents
\newpage

\section{Definitions}
We can calculate probabilities given some events already observed, said "conditioning on F" meaning "given F already observed".

Syntax:

\begin{empheq}[box=\fbox]{align}
    P(E|F) = \frac{P(EF)}{P(F)}
\end{empheq}

Reminders about general probability syntax:

\begin{itemize}
    \item $\cup$, OR, union
    \item $\cap$, AND, intersect. Shorthand $E,F$, $E \text{ and } F$, $EF$, $|EF|$
\end{itemize}

\begin{definition}[Chain Rule (aka Product Rule)]
    \begin{empheq}[box=\fbox]{align}
        P(EF) = \frac{P(E|F)}{P(F)}
    \end{empheq}
\end{definition}

The Chain Rule can be expanded with any number of events:

\begin{align}
    P(E_1,E_2,...,E_n) = P(E_1)P(E_2|E_1)...P(E_n|E_1,...,E_{n-1})
\end{align}

\begin{definition}[Law of Total Probability]
    \begin{empheq}[box=\fbox]{align}
        P(E) &= P(EF) + P(EF^C) \\
             &= P(E|F)P(F) + P(E|F^C)P(F^C)
    \end{empheq}
\end{definition}

\begin{definition}[Bayes Theorem]
    \begin{empheq}[box=\fbox]{align}
        P(F|E) = \frac{P(E|F)P(F)}{P(E)}
    \end{empheq}
\end{definition}

Remember that we can expand probability, we often need to use the expanded form for Bayes' theorem, using the Law of Total Probability for the denominator:

\begin{align}
    P(F|E) = \frac{P(E|F)P(F)}{P(E|F)P(F)+P(E|F^C)P(F^C)}
\end{align}

\end{document}